\begin{abstract}
Tool-using agents must authorize structured calls while reading untrusted
documents, webpages, and tool outputs, so indirect prompt injection can attach
attacker-chosen effects to otherwise legitimate work. We study the
agent-security boundary that precedes this decision: authorization logic can
enforce only the distinctions its observation interface preserves; if that
interface merges policy-distinct executions, every downstream authorizer must
admit an unauthorized execution or withhold authorized work. Defaults, aliases,
list expansion, and pre-state make this interface hard to construct.
Registration turns this boundary into an executable security contract: it
treats candidate descriptors as untrusted hypotheses, executes controlled
interventions in copied sandboxes, and asks a policy-separation oracle whether
an observed difference requires another decision. A merged policy-separating
pair is a failure certificate; counterexamples drive refinement, and each
frozen descriptor is bound to its implementation, interventions, and policy
family.

Across 264 post-freeze contexts for 11 tools from three public MCP
implementations, source execution leaves policy-mixed cells in tool-name,
raw-call, and common-field views, while a typed-effect candidate and a strong
state-aware request preserve every tested distinction. In a 1,024-context
authority protocol, both interfaces support exact decisions from one shared
ACL, capability, and delegation engine, showing that the method is
representation-agnostic. A pre-commit case study consumes a frozen interface
without placing authority in the LLM. The result is a falsification-first
method for auditing pre-commit authorization interfaces in tool-using agents.
\end{abstract}
